\section{A shrinking-core gradient and frequency barrier for unforced corrections}
\label{sec:xxxiii-shrinking}
In the authenticated forced construction of OpenAI \cite{openai2026},
Section~10, equations (10.20)--(10.21), along
$x_\tau=(\sqrt{2X_{\rm in}\tau},0,0)$,
\[
u_{B,\theta}(x_\tau,1-\tau)
=\tau^{-A}(e_0+O(\tau^{2h})),
\quad A=1/2+h,\quad 0<h<1/100,\quad e_0>0.
\]
The background also obeys a uniform whole-space kinetic-energy
bound $B=\sup_{0\le t<1}\|u_B(t)\|_2<\infty$.
These statements follow from the published Section~10 source,
without inventing a numerical cutoff profile.

\begin{lemma}[Two explicit three-dimensional interpolation bounds]
\label{lem:xxxiii-agn}
For a smooth vector field $v$ on $\mathbb R^3$ with
$v\in L^2$ and $\nabla v\in L^\infty$,
\[
\boxed{\quad\|v\|_\infty^5\le
\frac{24}{\pi}\|v\|_2^2\|\nabla v\|_\infty^3.\quad}
\]
For $v\in L^2$ with $\Delta v\in L^2$,
\[
\boxed{\quad\|v\|_\infty\le
C_2\|v\|_2^{1/4}\|\Delta v\|_2^{3/4},
\quad C_2=\frac4{\pi\sqrt2\,3^{7/8}}.\quad}
\]
Here $\|\nabla v\|_\infty$ denotes the operator-norm Lipschitz
constant, and $C_2$ is an admissible explicit constant,
not a claim of best possible interpolation constant.
\end{lemma}
\begin{proof}
At a point $|v(x_0)|=a>0$, the Lipschitz estimate gives
$|v|\ge a/2$ on $B(x_0,a/(2L))$, where
$L=\|\nabla v\|_\infty$. Integrating the squared norm
gives $\|v\|_2^2\ge\pi a^5/(24L^3)$ and then the
first bound by taking suprema.
For the second, use the unitary Fourier transform,
split at $|\xi|=K$, and use Cauchy--Schwarz:
\[
\|v\|_\infty\le
\frac{K^{3/2}}{\pi\sqrt6}\|v\|_2+
\frac{K^{-1/2}}{\pi\sqrt2}\|\Delta v\|_2.
\]
The exact radial integrals are
$\int_{|\xi|\le K}1\,d\xi=4\pi K^3/3$ and
$\int_{|\xi|>K}|\xi|^{-4}d\xi=4\pi/K$.
Choose $K^2=\|\Delta v\|_2/(\sqrt3\|v\|_2)$;
the two terms sum to the stated constant.
\end{proof}

\begin{theorem}[Necessary cancellation cost at the authentic singular trajectory]
\label{thm:xxxiii-cancellation}
Let $u$ be any smooth finite-energy, exactly unforced
three-dimensional Navier--Stokes comparison solution on
$[0,1)$ with $\|u(0)\|_2=x<\infty$.
Write $w=u-u_B$, $M=x+B$. The energy inequality and
triangle inequality give $\|w(t)\|_2\le M$ for all $t<1$.
Suppose that for all sufficiently small $\tau$ the
completed azimuthal velocity is no greater than
$(e_0/4)\tau^{-A}$ along the source path:
\[
|u_\theta(x_\tau,1-\tau)|\le
(e_0/4)\tau^{-A}.
\]
Then both the following lower bounds are \emph{necessary}:
\begin{align}
\|\nabla w(1-\tau)\|_\infty
&\ge\left(\frac{\pi(e_0/4)^5}{24M^2}\right)^{1/3}
\tau^{-5A/3},
\label{eq:xxxiii-gradient-cost}\\
\|\Delta w(1-\tau)\|_2
&\ge\left(\frac{e_0}{4C_2M^{1/4}}\right)^{4/3}
\tau^{-4A/3}.
\label{eq:xxxiii-h2-cost}
\end{align}
In particular, either reversed upper bound with the
same coefficient suffices to preserve a divergent
completed azimuthal velocity along $x_\tau$.
\end{theorem}
\begin{proof}
The source (10.21) is at least $(e_0/2)\tau^{-A}$
for all sufficiently small $\tau$. By the assumed
cancellation, $|w(x_\tau,1-\tau)|\ge
|w_\theta(x_\tau,1-\tau)|\ge(e_0/4)\tau^{-A}$.
Apply Lemma~\ref{lem:xxxiii-agn} to the correction,
using its energy bound $M$, to obtain both inequalities.
If either upper bound holds, the lemma instead bounds
$|w|\le(e_0/4)\tau^{-A}$ and the source's velocity
growth survives by the reverse triangle inequality.
\end{proof}

\begin{theorem}[Necessary high-frequency tail for cancellation]
\label{thm:xxxiii-tail}
With the same notation, define a moving Fourier cutoff
\[
K_\tau=\left(\frac{\pi\sqrt6\,e_0}{8M}\right)^{2/3}
\tau^{-2A/3},\qquad
P_{>K_\tau}=\mathbf 1_{\{|\xi|>K_\tau\}}(D).
\]
Under the cancellation assumption of
Theorem~\ref{thm:xxxiii-cancellation}, the correction
necessarily satisfies
\[
\boxed{\quad
|(P_{>K_\tau}w)(x_\tau,1-\tau)|
\ge(e_0/8)\tau^{-A}.\quad}
\]
Hence $|(P_{>K_\tau}w)(x_\tau,1-\tau)|=o(\tau^{-A})$
is a separate, sufficient spectral criterion for
preservation of the source's singularity.
The threshold scales as
$K_\tau\asymp\tau^{-1/3-2h/3}$, strictly slower
than the inverse radial similarity scale $\tau^{-1/2}$.
\end{theorem}
\begin{proof}
Unitary Fourier inversion and the volume of a ball give
$\|P_{\le K}w\|_\infty\le
MK^{3/2}/(\pi\sqrt6)$.
At $K=K_\tau$ this upper bound equals
$(e_0/8)\tau^{-A}$. The correction must have magnitude
at least $(e_0/4)\tau^{-A}$ at the source point;
the high-frequency piece therefore contributes at least
$(e_0/8)\tau^{-A}$. The exponent comparison follows
from $2A/3=1/3+2h/3<1/2$.
\end{proof}

\paragraph{Scope and release boundary.}
The source-specific ingredients used here are only the
actual Section~10 uniform energy and asymptotic growth,
not unverified numeric A.11/A.19 parameters.
The two Fourier constants, interpolation optimizer,
Gaussian solenoidal positive control, and 30,000 rational
dichotomy fixtures were independently reproduced.
These are necessary cancellation costs and sufficient
conditional retention inequalities, not an existence
proof for a nonzero-data unforced correction.
The exact correction $w=-u_B$ corresponding to the
trivial unforced solution $u=0$ is not excluded and
may satisfy the compulsory growth costs.
The full unforced Navier--Stokes singularity theorem
and kernel-checked Lean PDE proof remain open.
